% =====================================================================
% ROZDZIAŁ 7 — Zakończenie
% =====================================================================
\chapter{Zakończenie}
\label{rozd:zakonczenie}

\section{Odpowiedzi na pytania badawcze}
\label{sec:z-rq}

\paragraph{RQ1 (sygnał kaskadowy).} Tak: przejęte konto zgodne
z każdą statyczną relacją jest wykrywalne z dynamiki kaskady.
Hierarchia detektorów jest jednoznaczna (1-hop $0{,}70$; GNN
statyczny $0{,}76$; ręczny kontekst $0{,}81$; temporalny GNN
$0{,}87$ z czułością operacyjną $0{,}099$@1\%), przewaga jest
temporalna w sensie ablacyjnym (tasowanie: $0{,}87\to0{,}73$)
i przenosi się kierunkowo na trzy realne grafy (Enron: temporalny
$0{,}85$, czułość $0{,}215$ pod atakiem skrytym, przy czułości
$0{,}000$ detektora wolumetrycznego; SNAP: kontrast reżimowy
replikowany w kierunku na obu grafach).

\paragraph{RQ2 (granice).} Sygnał wymaga fan-outu $\ge2$ (atak
na jednego odbiorcę: AUC $\approx0{,}5$ dla każdego detektora);
propagacja nie jest konieczna (przy $p_{\mathrm{inf}}=0$ i fan-oucie
$\ge2$ wykrywalność wraca). Sygnał degraduje łagodnie ze skalą,
kurczy się pod szumem etykiet (nauczony model najszybciej:
$0{,}72\to0{,}59$ przy $20\%$), a między organizacjami najlepiej
transferują się cechy wolumetryczne. Przy realnej prewalencji
detektory są instrumentem triażu: $\sim\!10^3$ fałszywych alarmów
na wykrycie.

\paragraph{RQ3 (sygnał międzywarstwowy).} Tak: niespójność
międzywarstwowa oddziela impersonację zewnętrzną od legalnych
nowych nadawców (AUC $0{,}70\to0{,}88$, fałszywe alarmy
$100\%\to0\%$), kontrola tasowania spada do losowego, a kontrola
tylko-OSINT i tc-split lokalizują przyczynowy margines warstw
($+0{,}03$ AUC; $0{,}81\to0{,}90$ na plasterku nowych nadawców).
Uczciwa ablacja pokazuje, że większość sygnału rankingowego niesie
spójność czasowa, a wartość OSINT leży po stronie fałszywych
alarmów; sygnał eroduje pod realistyczną fabrykacją zgodnie ze
stwierdzeniem~\ref{prop:fab}.

\paragraph{RQ4 (narażenie i granice).} Narażenie napędza kanał
ataku, nie ranga ofiar: zimna impersonacja zostawia strukturalną
podłogę $\mathrm{Susc}\approx0{,}96$ na każdym stanowisku, a w kanale
lateral executive są \emph{najbardziej} narażeni --- bogactwo
osadzenia nie kupuje ochrony. Treść i graf zawodzą na różnych
atakach, a ich fuzja trzyma się blisko lepszego w każdym reżimie.
Najważniejsze odkrycie metodyczne rozprawy --- pułapka wycieku
rytmu --- pokazuje, że rytm estymowany na pełnej historii zawyża
czujność behawioralną z $0{,}00$ (uczciwie) do $0{,}73$ (z
wyciekiem): miara ceny naruszenia warunku przyczynowości
na realnym korpusie.

\section{Wkład rozprawy}
\label{sec:z-wklad}

Rozprawa wnosi: (i)~ramę formalną ewaluacji strumieniowej ---
klasę cech i detektorów zdarzeniowo-przyczynowych, cztery
mechanizmy wycieku jako twierdzenia z konstrukcyjnymi
kontrprzykładami, dekompozycję wyniku detektora na składową
czysto czasową i sprzęgającą czas z grafem, wynik o
nierozróżnialności śladów przy realistycznej fabrykacji
(rozdz.~\ref{rozd:formalny}); (ii)~pomiar i formalizację lateral
phishingu jako naruszenia grafu komunikacji w czasie --- z
odbiorczą falą kaskady jako niezmiennikiem, atakującym skrytym
jako testem krytycznym i replikacją na trzech realnych
topologiach (rozdz.~\ref{rozd:kaskada}); (iii)~pomiar
i formalizację impersonacji zewnętrznej jako naruszenia grafu
wiedzy domenowej --- multipleks, niespójność międzywarstwowa,
narażenie detektorowo-relatywne i granica graf--treść
(rozdz.~\ref{rozd:multipleks}); (iv)~metodologię ewaluacji
leak-aware, scalającą protokoły literaturowe z warunkami
twierdzonymi i operacjonalizowaną w paczce reprodukcyjnej, z
jawną arytmetyką bazy i granicami wypowiedzianymi
(rozdz.~\ref{rozd:synteza}).

\section{Ograniczenia}
\label{sec:z-ograniczenia}

Złośliwość jest syntetyczna: żaden publiczny korpus nie łączy
realnych kont z etykietowanymi przejęciami i ich propagacją, więc
walidujemy mechanizm (rozłożona w czasie kaskada separuje się
od ruchu legalnego dynamiką odbiorczą), nie detekcję realnych
intruzów; warstwy OSINT są syntetyczne także na realnym grafie.
Parametry kaskady i pacing są nasze; swepy pokazują stabilność
hierarchii, ale atak single-recipient pozostaje poza sygnałem
z konstrukcji. Skala jest ograniczona ($N{=}150$ w studium;
realne grafy do $\sim\!2$k węzłów), a graf studyjny jest gęstszy
niż typowa poczta korporacyjna --- nośnikiem mechanizmu poza
tym ustawieniem są replikacje Enron/SNAP, z zastrzeżeniem
kierunkowym (3--8 ziaren). Próg operacyjny pozostaje otwarty:
nawet najlepszy punkt pracy jest rzędy wielkości od progu
produkcyjnego ($<\!4$ FP/M); raportujemy tę lukę jako wymaganie
otwarte, nie rozwiązane.

\section{Kierunki przyszłej pracy}
\label{sec:z-przyszla}

Najbliższy krok to walidacja na korpusie z etykietowanymi,
\emph{realnymi} przejęciami --- zamiana dowodu mechanizmu
w pomiar wykrywalności. Model zagrożeń wymaga presji
adaptacyjnej: atakujący, który \emph{uczy się} grafu komunikacji
i dostosowuje pacing oraz wzorzec propagacji, by zejść pod falę
odbiorczą, wraz z dryfem populacji \citep{pendlebury2019tesseract}.
Po stronie modelowania: bogatszy kontekst relacyjny --- hierarchia,
kompetencje i publiczny ślad jako warstwy na tych samych kontach ---
i synteza obu kanałów w jednym detektorze fuzji graf--treść.
Po stronie ewaluacji: rozdział formalny sugeruje naturalne
rozszerzenia --- kontrola tasowania per kubełek (izolująca
czystą składową porządku od ko-występowania) oraz fabrykator
optymalizujący spójność, nie istnienie śladów. Wynik odbiorczy
jest naturalnym wejściem do kolejki triażu SOC --- pod warunkiem
ewaluacji, która wiąże go z realną redukcją obciążenia analityka
przy stałym budżecie alertów.
